\section{Longitudinal regime contrasts: preserving the effective uncertainty method}
\label{exec23:ltmle-regimes}

A servicing proposal can compare a fixed contact schedule with a rule that
responds to hardship observed after an earlier contact. The distinction matters
because the first contact may change both hardship and the second action.
The selected ltmle source provides an R reference for binary treatment regimes,
longitudinal treatment and censoring histories, and treatment-specific means.
Its \texttt{ltmle} interface handles a single regime or a pair; its
\texttt{ltmleMSM} interface also supports a working marginal structural model.
A general working model is a projection under specified weights and covariates;
it need not reproduce each separate regime mean. A saturated two-regime
parameterization is a special case. The number of policies alone does not
determine whether a working model is saturated.\citep{exec23_ltmle_source}

The inspected revision declares version 1.3-0 and GPL-2. The main source,
package declaration and supporting files match the complete tree of commit
\texttt{338c029dae9692ef20714125773da7037688993b}, dated 17 April 2023.
A new default-branch query restricted to the client cutoff returns this same
commit. Repository metadata also records a June 2026 push; that event alone
does not establish a newer default-branch estimator or equality with a CRAN
artifact. The declared imports are \texttt{Matrix} and \texttt{matrixStats};
\texttt{SuperLearner} is suggested and required by the inspected non-GLM path.
The default \texttt{SL.library="glm"} uses generalized linear models. It does
not request an ensemble merely because an ensemble option exists.
\citep{exec23_ltmle_source}

\subsection{A dynamic intervention and its information clock}

Let $H$ indicate baseline borrower risk, with $P(H=1)=1/3$. The first contact
is $A_0$, hardship measured after it is $L_1$, and the second contact is $A_1$.
The outcome $Y=1$ denotes default by the fixed horizon. The information order
is $H,A_0,L_1,A_1,C,Y$, where $C=1$ means the final outcome is observed.
For an exact authored mechanism, stipulate
\begin{align}
 P(L_1=1\mid H,A_0)
 &=\begin{cases}
 1/4,&(H,A_0)=(0,0),\\
 1/8,&(H,A_0)=(0,1),\\
 3/4,&(H,A_0)=(1,0),\\
 1/2,&(H,A_0)=(1,1),
 \end{cases}\label{exec23:ltmle-hardship}\\
 P(Y=1\mid H,A_0,L_1,A_1)
 &=\frac1{10}+\frac H{10}+\frac{L_1}{5}
       -\frac{A_0}{25}-\frac{A_1}{20}.
 \label{exec23:ltmle-default}
\end{align}
Every binary history has a valid default probability: the minimum is $1/100$
and the maximum is $2/5$. Conditional on the listed history, let $C$ be
independent of $Y$, with probability $9/10$ of observation when $L_1=0$ and
$3/4$ when $L_1=1$. This is a stipulated observation mechanism; absence of
outcome data is not redefined as absence of default.

Compare status quo $d^0=(0,0)$, targeted servicing
$d^T=(H,L_1)$, and broad contact $d^B=(1,1)$. The second targeted action reads
the hardship observed at its decision time, including any effect of the first
action on hardship. To estimate these targets from logged data, one would need
consistency, sequential exchangeability and treatment support, together with
the censoring assumptions and relevant estimator regularity. These assumptions
are separate from the mathematical truth of the stipulated mechanism.
Its regime mean is
\begin{equation}
 \theta_d=\sum_{h,l}P(H=h)P(L_1=l\mid h,A_0=d_0(h))
 E[Y\mid h,A_0=d_0(h),l,A_1=d_1(l)].
 \label{exec23:ltmle-gformula}
\end{equation}
The hardship means under the three policies are $5/12,1/3,1/4$.
Thus both path enumeration and the separate expectation identity
$\theta_d=1/10+1/30+E_d(L_1)/5-E_d(A_0)/25-E_d(A_1)/20$
give the same risks. If a default costs USD 1,000 and contacts cost USD 6
and USD 4 at the two times, respectively, define net value relative to status quo by
\begin{equation}
 V_d=1000(\theta_{d^0}-\theta_d)-6E_d(A_0)-4E_d(A_1).
 \label{exec23:ltmle-value}
\end{equation}
The default cost is illustrative and fixed; it excludes any unmodeled changes
in recovery, customer behavior or subsequent service cost.

\begin{center}\small
\begin{tabular}{lccc}
\toprule
Policy & Default risk & Contact cost (USD) & Net value (USD)\\
\midrule
Status quo & $13/60$ & $0$ & $0$\\
Targeted dynamic & $17/100$ & $10/3$ & $130/3$\\
Broad contact & $7/75$ & $10$ & $340/3$\\
\bottomrule
\end{tabular}\end{center}
These are exact per-account expectations under the assumed law, not estimates
returned by ltmle or evidence of bank benefit. The positive broad-contact value
depends on the assumed effect and cost parameters. The example gives a target
for a future package comparison; it does not select a bank policy.
\citep{exec23_ltmle_exact}

Censoring alters even this small example if the analyst drops missing outcomes.
The complete-case risk under status quo is $69/335$, compared with the target
$13/60$; its bias is $-43/4020$. Under targeted servicing, the corresponding
risks are $139/850$ and $17/100$, giving bias $-11/1700$. Write
$\pi(L_1)=P(C=1\mid L_1)$. Under the stipulated conditional independence,
\begin{equation}
 E_d\!\left\{\frac{CY}{\pi(L_1)}\right\}
 =E_d\!\left[E\!\left\{\frac{CY}{\pi(L_1)}\,
          \middle|\,H,A_0,L_1,A_1\right\}\right]
 =E_d(Y)=\theta_d.
 \label{exec23:ltmle-censoring}
\end{equation}
Exhaustive enumeration verifies this population identity for all three
policies and checks that their probability masses sum to one. Known
observation probabilities make this an algebraic check. Estimating those
probabilities, multiplying treatment probabilities, and obtaining useful
uncertainty require additional work. The bound $\pi\geq3/4$ in this example
does not demonstrate overlap in a real servicing log.
\citep{exec23_ltmle_exact}

\subsection{Why a named variance option is insufficient}

The source defaults to \texttt{variance.method="tmle"}, but
\texttt{CreateInputs} switches to \texttt{ic} when the additional variance
method is unavailable. In the inspected revision, non-null \texttt{id},
nonbinary outcomes, \texttt{gcomp}, stratification and a deterministic-Q
function each make that additional method unavailable. The fallback itself
does not always warn: a separate warning path checks selected positivity and
rare-event indicators. The absence of that warning is therefore not evidence
that the caller's requested method ran.\citep{exec23_ltmle_source}

For a single-regime TMLE summary, the code takes the larger of the available
influence-curve and additional variance estimates. For paired regimes, it
applies each contrast gradient to the covariance matrices and selects the
larger available standard deviation for that reported measure. It does not
justify a new covariance matrix formed by taking every element's maximum.
Risk and odds ratios are removed for nonbinary outcomes. Transformations back
to the outcome scale, including the distinction between mean and difference,
must be preserved. A larger variance estimate alone is not a proof of
coverage.\citep{exec23_ltmle_source}

For clustered records, \texttt{HouseholdIC} sums record contributions within
each identifier and multiplies them by $m/n$, where $m$ is the number of
identifiers and $n$ is the number of records. Subsequent summary calculations
use the number of influence rows. This keeps the record-weighted target;
it does not automatically change that target to an equally weighted mean
of households of different sizes. The choice of weights and scientific
sampling unit must precede the conversion.

Consider a separate exact covariance fixture with four household influence
pairs $(-3,-2),(-1,0),(1,0),(3,2)$, each repeated twice as records. These are
hypothetical influence values, not fitted ltmle output. Aggregation and
rescaling recover the four pairs. Dividing their sample covariance by four
gives the covariance of the two means as
\begin{equation}
 \widehat\Sigma=\begin{pmatrix}5/3&1\\1&2/3\end{pmatrix},\qquad
 \widehat{\operatorname{Var}}(\widehat\theta_1-\widehat\theta_0)
 =5/3+2/3-2=1/3.
 \label{exec23:ltmle-cluster}
\end{equation}
Dropping the covariance instead gives $7/3$. Treating the eight duplicated
records as independent gives variance $5/7$ for the first mean, instead of
$5/3$. Their exact ratio is $7/3$ because sample variance uses different
degrees of freedom. The often-quoted factor of two is the leading
population-variance comparison for two identical records per independent
household. Neither the four-household calculation nor a small-sample
$t$ critical value proves calibration. ltmle uses a $t$ critical value below
100 influence rows and a normal critical value otherwise; the sampling law
and nuisance remainder still need justification.
\citep{exec23_ltmle_source,exec23_ltmle_exact}

\subsection{A proposed R comparator and the alternatives}

A candidate interface should declare the node order, policy information set,
binary action coding, event/censoring coding, outcome range, target weights
and cluster identity. Dynamic rules must not read future outcomes merely
because the R function receives a complete row. Numeric censoring inputs
use zero for censored and one for uncensored; the named conversion helper
distinguishes that meaning from a censoring indicator. For absorbing events,
the source checks that later outcome nodes remain one before cleaning
post-event records. The shorthand that later data are ``ignored'' must not
erase that validation requirement.\citep{exec23_ltmle_source}

The cumulative-probability floor defaults to $0.01$ after multiplying the
incremental factors. It is not a floor on each factor before multiplication.
The source returns bounded and unbounded cumulative probabilities, which
support diagnostics; neither operation proves positivity. The outcome-range
code can truncate values before scaling and has a division by the range
width, motivating an explicit positive-width contract for a future adapter.
No warning should be interpreted as a complete fit audit: inspected helpers
suppress some rank-deficient prediction and extreme-probability warnings,
and the targeting helper can suppress a nonconvergence warning before its
score-repair path. A failed matrix solve can return missing standard errors.
Unavailable uncertainty must remain an explicit result state.
\citep{exec23_ltmle_source}

The proposal retains this version as a technical reference and a candidate
native R comparator. Direct integration requires resolved dependencies,
distribution rights, observed fit and warning behavior, and matched tests
against the declared policy and covariance targets. The GPL-2 declaration
and license distinguish running software from distributing or modifying it;
a separate process is not, by itself, a legal compatibility determination.
An integration design must review the actual distribution and dependency
chain. Porting to Python, Java or C++ additionally requires parity for
history/rule construction, censoring, outcome scaling, cumulative bounds,
learner behavior, clustering and covariance. No measured deployment or
performance need currently justifies that larger porting cost.
\citep{exec23_ltmle_source}

lmtp and gfoRmula remain substantive alternatives with their own modified
treatment, simulation and uncertainty contracts. Their names do not establish
equivalent interventions or interchangeable output. Selection should follow
the requested regime and observation law, then compare full fitting,
diagnostics, uncertainty, failure handling and maintenance costs. This
source-level reference decision grants no native integration capability.
The exact examples specify what a proposed trial should preserve; package
execution and empirical qualification remain distinct evidence.

